\chapter{Incontinence: An Overview}
\index{Incontinence: An Overview}

\section*{Definition and Overview}
Incontinence is the involuntary loss of urine or faeces, or both. Urinary incontinence covers stress, urge, mixed, and overflow types, while faecal incontinence is the inability to control the passage of stool or wind. The two frequently occur together, especially in older people, after childbirth, and in people with neurological conditions. Incontinence is not a disease in itself but a symptom with many possible causes, and it is never a normal or acceptable part of ageing. It affects physical health, dignity, independence, and quality of life, and it is one of the main reasons older people move into residential care. With modern assessment and treatment, most people can be substantially helped, and many can become fully continent.

\section*{Epidemiology}
Incontinence is very common and becomes more so with age. Around one in three women and one in ten men experience urinary incontinence, and around one in 20 adults experience faecal incontinence at some time. Rates are higher in older people, in women who have had vaginal births, and in people living in residential care, where around half of residents have some degree of incontinence. It is estimated that only about a third of affected people seek help. The financial cost to health systems, and the personal cost in lost work, social participation, and wellbeing, is substantial, and the burden is set to grow as the population ages.

\section*{Aetiology and Pathophysiology}
Continence relies on a coordinated system: muscles of the pelvic floor and sphincters hold urine and stool in, muscles of the bladder and bowel wall relax and contract to store and expel, and nerves link these structures to the brain. Incontinence results when any of these elements fails. Urinary stress incontinence follows pelvic floor weakness, urge incontinence follows overactivity of the bladder muscle, and overflow incontinence follows obstruction or a weak bladder muscle. Faecal incontinence can follow anal sphincter damage from childbirth, chronic constipation and straining, diarrhoea, rectal disease, or reduced awareness of the need to defecate. Neurological conditions including stroke, multiple sclerosis, Parkinson's disease, dementia, and spinal cord injury can affect either or both systems.

\section*{Clinical Features}
\begin{itemize}
    \item Leakage of urine with coughing, sneezing, laughing, or exercise
    \item Sudden urgent need to pass urine with leakage before reaching the toilet
    \item Frequent urination, including through the night
    \item Accidental passage of stool or inability to control wind
    \item Urgency to open the bowels, sometimes with leakage on the way to the toilet
    \item Staining of underclothes and the need to wear pads
    \item Skin soreness, rashes, and recurrent urinary or skin infections
    \item Avoidance of socialising, exercise, travel, and intimacy
    \item In older people, falls while hurrying to the toilet and increased dependence on carers
\end{itemize}

\section*{Differential Diagnosis}
The cause of incontinence is identified through careful assessment, because different types need different treatment.
\begin{itemize}
    \item \textbf{Transient causes:} urinary tract infection, constipation, medicines, excessive fluid intake, and delirium
    \item \textbf{Stress versus urge versus mixed urinary incontinence:} distinguished by history and examination
    \item \textbf{Overflow incontinence:} obstruction or poor bladder emptying
    \item \textbf{Faecal impaction:} a common and reversible cause of both faecal and urinary leakage
    \item \textbf{Neurological disease:} spinal cord problems, stroke, or neuropathy
    \item \textbf{Prostate disease in men:} enlargement, cancer, or treatment effects
    \item \textbf{Structural problems:} pelvic organ prolapse, fistulas, or anorectal disease
    \item \textbf{Functional problems:} poor mobility, dexterity, or cognition preventing toilet use
\end{itemize}

\section*{When to Seek Medical Care}
Anyone with leakage that affects daily life, dignity, or sleep should be assessed, since treatment works best when started early. Seek review for new or sudden onset incontinence, incontinence with pain or blood, or leakage accompanied by other neurological symptoms. Carers should raise concerns for people who cannot speak for themselves, particularly with sudden changes.

\textbf{Contact your local emergency services for:}
\begin{itemize}
    \item Sudden inability to pass urine with severe pain, indicating acute urinary retention
    \item New loss of bladder or bowel control with back pain, leg weakness, or numbness, which may indicate spinal cord compression
    \item Sudden neurological symptoms such as face, arm, or leg weakness or difficulty speaking
\end{itemize}

\section*{Diagnosis and Investigations}
\begin{itemize}
    \item \textbf{Detailed history:} type, timing, and pattern of leakage, fluid and diet intake, medicines, and obstetric or surgical history
    \item \textbf{Bladder and bowel diaries:} several days of recording intake, output, and leakage events
    \item \textbf{Urinalysis:} to exclude infection, blood, or glucose
    \item \textbf{Physical examination:} abdomen, pelvic floor, and in women a vaginal examination; a rectal examination when faecal symptoms are present
    \item \textbf{Post-void residual:} ultrasound measurement of urine left in the bladder
    \item \textbf{Anorectal assessment:} examination, and in selected cases anorectal manometry or ultrasound
    \item \textbf{Urodynamics:} specialised bladder function testing where the diagnosis is unclear or surgery is planned
    \item \textbf{Referral:} to continence services, physiotherapy, urology, colorectal surgery, or aged-care specialists
\end{itemize}

\section*{Management}
\begin{itemize}
    \item \textbf{Pelvic floor muscle training:} first-line for stress and mixed urinary incontinence
    \item \textbf{Bladder training and scheduled voiding:} for urgency and urge incontinence
    \item \textbf{Bowel training and diet:} regular toilet habits, adequate fibre and fluid, and treatment of constipation
    \item \textbf{Lifestyle changes:} weight loss, stopping smoking, and reducing caffeine and alcohol
    \item \textbf{Medicines:} drugs for an overactive bladder, laxatives or antidiarrhoeals, and vaginal oestrogen where appropriate
    \item \textbf{Procedures and surgery:} sling procedures, botulinum toxin, sacral nerve stimulation, or sphincter repair for selected cases
    \item \textbf{Products and aids:} pads, protective garments, catheter products, and toileting aids that maintain dignity and independence
    \item \textbf{Continence products funded by schemes:} access to subsidised products through state and federal programs
    \item \textbf{Carer support and care planning:} involving family, aged-care services, and allied health professionals
\end{itemize}

\section*{Complications}
\begin{itemize}
    \item Skin breakdown, pressure injuries, and fungal or bacterial skin infections
    \item Recurrent urinary tract infections
    \item Falls and fractures, especially at night
    \item Sleep disturbance and chronic fatigue
    \item Depression, anxiety, social withdrawal, and loss of self-esteem
    \item Strained relationships and reduced intimacy
    \item Loss of employment, independence, and community participation
    \item Admission to residential care that might otherwise have been avoidable
\end{itemize}

\section*{Prognosis and Outcomes}
The outlook depends on the type and cause, but most people improve with structured management. Pelvic floor training, bladder training, medicines, and lifestyle change help the majority of people with urinary incontinence, and many are cured. Faecal incontinence from reversible causes such as constipation or diarrhoea commonly resolves with treatment, and sphincter repairs help many women with childbirth-related damage. Even when full continence is not achievable, modern products and care plans allow people to live active, dignified lives. Continence assessment should be offered proactively, because untreated incontinence drives so many avoidable complications.

\section*{Prevention and Self-Care}
\begin{itemize}
    \item Do pelvic floor exercises regularly, especially after childbirth
    \item Maintain a healthy weight and stay physically active
    \item Eat a high-fibre diet and drink enough fluid to avoid constipation
    \item Go to the toilet when needed rather than holding on for long periods
    \item Limit caffeine and alcohol, and avoid excessive evening fluids
    \item Stop smoking to reduce coughing and straining
    \item Seek help early from a doctor or continence service
    \item Use appropriate pads and skin-care products to protect the skin
    \item If caring for someone, respond promptly to toilet requests and maintain their dignity
    \item Raise incontinence concerns openly in health checks, as it is treatable
\end{itemize}

\section*{Sources and Further Reading}
The following reputable sources provide further information for healthcare students and patients seeking reliable, current advice about incontinence.
\begin{itemize}
    \item Continence Foundation of Australia. \textit{Bladder and bowel health information.} https://www.continence.org.au/
    \item Healthdirect Australia. \textit{Incontinence overview.} https://www.healthdirect.gov.au/
    \item National Continence Program. \textit{Government support and products.} https://www.health.gov.au/
    \item NHS. \textit{Urinary incontinence and bowel control.} https://www.nhs.uk/
    \item Mayo Clinic. \textit{Incontinence.} https://www.mayoclinic.org/
    \item International Continence Society. \textit{Patient resources.} https://www.ics.org/
\end{itemize}
